% $Header: /cvsroot/latex-beamer/latex-beamer/solutions/generic-talks/generic-ornate-15min-45min.en.tex,v 1.4 2004/10/07 20:53:08 tantau Exp $

%\documentclass[trans]{beamer}
%\documentclass[draft]{beamer}
\documentclass{beamer}

\usepackage{enumerate}
\usepackage{amscd,amsmath,xspace,amssymb,amsfonts,mathrsfs,latexsym}
\usepackage{graphicx}

\usepackage{array}
%\usepackage{mbboard}
%\usepackage{multicol}
%\usepackage[pdftex]{color}

%\includeonlyframes{current}

\usepackage[all,poly,knot,tips]{xy}

\newcommand{\lra}{\longrightarrow}
\newcommand{\lla}{\longleftarrow}
\newcommand{\Lra}{\Longrightarrow}
\newcommand{\Lla}{\Longleftarrow}

%% Script Shortcuts 
\def\CA{{\mathscr A}}
\def\CB{{\mathscr B}}
\def\CC{{\mathscr C}}
\def\CD{{\mathscr D}}
\def\CE{{\mathscr E}}
\def\CF{{\mathscr F}}
\def\CG{{\mathscr G}}
\def\CH{{\mathscr H}}
\def\CI{{\mathscr I}}
\def\CJ{{\mathscr J}}
\def\CK{{\mathscr K}}
\def\CL{{\mathscr L}}
\def\CM{{\mathscr M}}
\def\CN{{\mathscr N}}
\def\CO{{\mathscr O}}
\def\CP{{\mathscr P}}
\def\CQ{{\mathscr Q}}
\def\CR{{\mathscr R}}
\def\CS{{\mathscr S}}
\def\CT{{\mathscr T}}
\def\CU{{\mathscr U}}
\def\CV{{\mathscr V}}
\def\CW{{\mathscr W}}
\def\CX{{\mathscr X}}
\def\CY{{\mathscr Y}}
\def\CZ{{\mathscr Z}}

\newcommand{\A}{\ensuremath{\mathcal{A}}\xspace}
\newcommand{\B}{\ensuremath{\mathcal{B}}\xspace}
\newcommand{\Af}{\ensuremath{\mathcal{A}_f}\xspace}
\newcommand{\C}{\ensuremath{\mathcal{C}}\xspace}
\newcommand{\D}{\ensuremath{\mathcal{D}}\xspace}
\newcommand{\F}{\ensuremath{\mathcal{F}}\xspace}
\newcommand{\I}{\ensuremath{\mathcal{I}}\xspace}
\newcommand{\K}{\ensuremath{\mathcal{K}}\xspace}
\newcommand{\Kt}{\ensuremath{\mathcal{K}_t}\xspace}
\newcommand{\Kl}{\ensuremath{\mathcal{K}_{\ell}}\xspace}
\renewcommand{\L}{\ensuremath{\mathcal{L}}\xspace}
\renewcommand{\S}{\ensuremath{\mathcal{S}}\xspace}
\newcommand{\T}{\ensuremath{\mathcal{T}}\xspace}
\newcommand{\V}{\ensuremath{\mathcal{V}}\xspace}
\newcommand{\Vf}{\ensuremath{\mathcal{V}_f}\xspace}
\newcommand{\W}{\ensuremath{\mathcal{W}}\xspace}
\newcommand{\X}{\ensuremath{\mathcal{X}}\xspace}

\newcommand{\op}{\ensuremath{^{\textrm{op}}}\xspace}
\newcommand{\ps}{\ensuremath{^{\textrm{ps}}}\xspace}
\newcommand{\lax}{\ensuremath{^{\textrm{lax}}}\xspace}
\newcommand{\oplax}{\ensuremath{^{\textrm{oplax}}}\xspace}


\newcommand{\EMK}{\textnormal{\sf EM(\K)}\xspace}
\newcommand{\EM}{\textnormal{\sf EM}\xspace}
\newcommand{\EMwK}{\ensuremath{\EM^w(\K)}\xspace}

%\newcommand{\Cat}{\textnormal{\bf Cat}\xspace}
\newcommand{\Lex}{\textnormal{\sf Lex}\xspace}
\newcommand{\Lan}{\textnormal{\sf Lan}\xspace}
\newcommand{\Hom}{\textnormal{\sf Hom}\xspace}
\renewcommand{\t}{\times}
\newcommand{\ot}{\otimes}
\newcommand{\two}{\mathbf{2}\xspace}


\newtheorem{remark}[theorem]{Remark}
\newtheorem{proposition}[theorem]{Proposition}

\beamerdefaultoverlayspecification{<+->}

%\newenvironment{\next}{\uncover<+->{}{}}

\mode<presentation>
{
%  \usetheme[compress,subsection=false]{Szeged}
  % or ...
%  \useoutertheme[subsection=false]{miniframes}
%\useoutertheme[subsection=false, footline=authortitle, frame number]{miniframes}
\useoutertheme{infolines}
\setbeamercolor{separation line}{use=structure,bg=structure.fg!50!bg}
\usecolortheme{whale}
\usecolortheme{orchid}
\usecolortheme{rose}
\setbeamercolor{titlelike}{parent=structure}


%  \setbeamercovered{transparent}
  % or whatever (possibly just delete it)
}


\usepackage[english]{babel}
% or whatever

\usepackage[latin1]{inputenc}
% or whatever

%\usepackage{times}
%\usepackage[T1]{fontenc}
% Or whatever. Note that the encoding and the font should match. If T1
% does not look nice, try deleting the line with the fontenc.


\title%[Short Paper Title] % (optional, use only with long paper titles)
{\textbf{Skew monoidal categories and bialgebroids}}

%{Presentation Subtitle} % (optional)

\author{Ross Street}
\institute{Macquarie University}
%[Universities of Somewhere and Elsewhere] % (optional, but mostly needed)
%{
%  \inst{1}%
%  Department of Computer Science\\
%  University of Somewhere
%  \and
%  \inst{2}%
%  Department of Theoretical Philosophy\\
%  University of Elsewhere}
% - Use the \inst command only if there are several affiliations.
% - Keep it simple, no one is interested in your street address.

\date[Sept-Oct 2013]{\small Algebra Special Session,
57th Annual Meeting of AustMS,
University of Sydney} 

\subject{Algebra Special Session}
% This is only inserted into the PDF information catalog. Can be left
% out. 

% Delete this, if you do not want the table of contents to pop up at
% the beginning of each subsection:
\AtBeginSubsection[]
{
  \begin{frame}<beamer>
    \frametitle{Outline}
    \tableofcontents[currentsection]
  \end{frame}
}
% If you wish to uncover everything in a step-wise fashion, uncomment
% the following command: 

\beamerdefaultoverlayspecification{<.->}

\begin{document}

%  \begin{frame}
%    \titlepage
%  \end{frame}


\begin{frame}%[plain]
%  \frametitle{titlepage}

  \begin{centering}
    {\usebeamerfont{title}\usebeamercolor[fg]{title}\inserttitle

\medskip 

{\small Ross Street}

{\small Macquarie University} 

\medskip
\uncover<1->{{}}
}

\vfill

\centerline{\usebeamerfont{title}\usebeamercolor[fg]{title}\insertdate}
% CT09, Cape Town}

  \end{centering}


\end{frame}


% \begin{frame}
%   \frametitle{Outline}
%   \tableofcontents
%   % You might wish to add the option [pausesections]
% \end{frame}

%\section[Stricts and weaks]{The strict/weak dichotomy}

%\subsection{Classical universal algebra}
\begin{frame}
  \frametitle{Outline}
\begin{itemize}
\item {\em Skew monoidal categories} and other kinds of lax monoidal categories were considered by category theorists over the years.
\pause
\item The subject and the name became active last year when Kornel Szlach\'anyi related them to the bialgebroids as studied in the quantum group literature.
\pause
\item The work in this talk is joint with Steve Lack.
\pause
\item I shall begin by defining {\em flocks} and giving examples.
\pause
\item Then we will look at skew monoidal categories and examples.
\pause 
\item Categories themselves and bialgebroids will be characterized as closed skew monoidal structures.
\pause 
\item We will conclude by discussing the free skew monoidal category and `coherence'.
\end{itemize}
\end{frame}

\begin{frame}
  \frametitle{Herds}

 \begin{itemize}
\item A herd is a group without a distinguished leader. 
\pause
\item More precisely, a {\em herd} $A$ is a set with a ternary operation 
$$q : A\times A \times A \lra A$$ 
satisfying $$q(a,b,q(c,d,e))=q(q(a,b,c),d,e) \ ,$$
$$q(a,b,b) = a \ \ \text{and} \ \  b = q(a,a,b)\ .$$
\pause
\item Each group $G$ is a herd with $q(a,b,c)=ab^{-1}c$.
\pause
\item For any herd, choice of any element $j\in A$ makes $A$ a group with
multiplication $ab = q(a,j,b)$, identity element $j$ and $a^{-1}=q(j,a,j)$.
\pause
\item Let $\mathrm{Iso}_{\CC}(X,Y)$ be the set of invertible morphisms 
$a : X\lra Y$ in any category $\CC$. It becomes a herd via $q(a,b,c)=ab^{-1}c$.

Choice of $j\in \mathrm{Iso}_{\CC}(X,Y)$ gives a group isomorphic to the group
$\mathrm{Aut}_{\CC}(X)$ of automorphisms of $X\in \CC$. 
\end{itemize}
\end{frame}

\begin{frame}
  \frametitle{Flocks}

 \begin{itemize}
\item Flocks categorify herds. 
\pause
\item More precisely, a {\em right flock} $\CA$ is a category with a functor 
$$Q : \CA\times \CA^{\mathrm{op}} \times \CA \lra \CA$$ 
equipped with natural transformations
\begin{gather*}
\phi_{ABCDE} :Q( A,B,Q(C,D,E)) \longrightarrow Q(Q(A,B,C),D,E)
\end{gather*}
\begin{gather*}
\xi _{A}^{B}:Q( A,B,B) \longrightarrow A
\end{gather*}
\begin{equation*}
\zeta _{B}^{A}:B\longrightarrow Q( A,A,B) \ ,
\end{equation*}
\end{itemize}
\end{frame}
\begin{frame}
such that the following five conditions hold:
\begin{gather*}
\begin{split}
\xymatrix{
Q(Q(A,B,C),D,Q(E,F,G)) \ar[d]_{\phi}="1" & Q(A,B,Q(C,D,Q(E,F,G))) \ar[l]_{\phi} \ar[dd]^{Q(1,1,\phi)}\\
Q(Q(Q(A,B,C),D,E),F,G)  &  \\
Q(Q(A,B,Q(C,D,E)),F,G)\ar[u]^{Q(\phi,1,1)}  & Q(A,B,Q(Q(C,D,E),F,G) \ar[l]^{\phi}
} 
\end{split}\\ 
\begin{split}
\xymatrix{
Q(A,B,Q(B,B,C)) \ar[rr]^-{\phi} && Q(Q(A,B,B),B,C) \ar[d]^{Q(\xi,1,1)}="2" \\
Q(A,B,C) \ar[u]^{Q(1,1,\zeta)}="1" \ar[rr]_{1} && Q(A,B,C)
} 
\end{split}
\end{gather*}
\end{frame}
\begin{frame}
\begin{equation*}
\xymatrix{
Q(A,A,Q(B,C,D)) \ar[rr]^{\phi}   && Q(Q(A,A,B),C,D)  \\
& Q(B,C,D) \ar[lu]^{\zeta} \ar[ru]_{Q(\zeta,1,1)}  &
}
\end{equation*}
\begin{equation*}
\xymatrix{
Q(A,B,Q(C,D,D)) \ar[rd]_{Q(1,1,\xi)}\ar[rr]^{\phi}   && Q(Q(A,B,C),D,D) \ar[ld]^{\xi} \\
& Q(A,B,C)  &
}
\end{equation*}
\begin{gather*}
\begin{split}
\xymatrix{
&& Q(A,A,A)\ar[rrd]^{\xi} && \\
A \ar[rru]^{\zeta} \ar[rrrr]_-{1} &&  &&  A \ . \ \square 
} 
\end{split}
\end{gather*}
\end{frame}

\begin{frame}
\begin{example}
Consider the category $\mathrm{Vect}$ of finite dimensional (right) vector spaces over some fixed field $k$.
For $V\in \mathrm{Vect}$, write $V^*$ for the vector space of linear functions $\ell : V\lra k$. 
We have linear functions $\varepsilon : V^*\otimes V \lra k$ and $\eta : k \lra V\otimes V^*$
defined by $\varepsilon (\ell \otimes v) =\ell (v)$ and $\eta (1) = \sum_{i=1}^n e_i\otimes e^*_i$
where $\sum_{i=1}^n e_ie^*_i(v)=v$.  
Put $$Q(U,V,W)=U\otimes V^*\otimes W \ .$$ 
Take $\phi$ to be the identity function of
$U\otimes V^*\otimes W\otimes X^*\otimes Y$. 
Take $$\xi = 1_U\otimes \varepsilon : U\otimes V^*\otimes V \lra U$$ and
$$\zeta = \eta \otimes 1_V : V \lra U\otimes U^*\otimes V \ .$$
The five flock axioms can be verified.
\end{example}

\pause
A flock with $\phi$ invertible is called {\em Hopf}.
\end{frame}
\begin{frame}
\begin{example}
We make the category $\mathrm{Ab}^R$ of left $R$-modules into a flock.
Put $$Q(L,M,N)=\mathrm{Hom}_R(\mathrm{Hom}(N,M),L) \ .$$
\pause
We have the canonical left R-module morphism
$$\tau : \mathrm{Hom}(N,M) \otimes_R \mathrm{Hom}(L,K) \lra \mathrm{Hom}(\mathrm{Hom}_R(\mathrm{Hom}(N,M),L),K)$$
defined by $\tau (f\otimes_Rg)(\theta)=g(\theta(f))$.
Applying the contravariant functor $\mathrm{Hom}_R(-,H)$ to $\tau$ 
and using the $\otimes_R$-$\mathrm{Hom}_R$ adjunction, we obtain our  
\begin{eqnarray*}
\phi : \mathrm{Hom}_R(\mathrm{Hom}(\mathrm{Hom}_R(\mathrm{Hom}(N,M),L),K),H) \\
\lra \mathrm{Hom}_R(\mathrm{Hom}(N,M),\mathrm{Hom}_R(\mathrm{Hom}(L,K),H)) \ .
\end{eqnarray*}
\pause
Moreover, $\xi :\mathrm{Hom}_R(\mathrm{Hom}(M,M),L)\lra L$ evaluates at $1_M$ 

and
$\zeta : L\lra \mathrm{Hom}_R(\mathrm{Hom}(M,L),L)$ is $\zeta (x)(\theta)=\theta (x)$.
\end{example}
\end{frame}

\begin{frame}
\begin{example}
Again we make the category $\mathrm{Ab}^R$ of left $R$-modules into a flock.

Put $$Q(L,M,N)=N\otimes \mathrm{Hom}_R(M,L) \ .$$
\pause
The obvious morphism
$$\mathrm{Hom}_R(M,L)\otimes \mathrm{Hom}_R(K,H) \lra \mathrm{Hom}_R(M,L\otimes \mathrm{Hom}_R(K,H)$$
yields $\phi$ on applying the covariant functor $N\otimes-$.

\pause
\bigskip
Take $\xi : M\otimes \mathrm{Hom}_R(M,L) \lra L$ to be evaluation 
\bigskip
and $\zeta : M \lra M\otimes \mathrm{Hom}_R(L,L)$ to be $\zeta(x) = x\otimes 1_L$. 

\end{example}
\end{frame}
\begin{frame}
  \frametitle{Skew monoidal categories}

A category $\CR$ is {\em right skew-monoidal} when it is equipped with 

  \begin{itemize}
\item<1-> a functor $\otimes : \CR \times \CR \lra \CR$ and object $I\in \CR$ 

\item<2-> natural families of (not necessarily invertible) morphisms 
$$\alpha : A\otimes (B\otimes C) \lra (A\otimes B)\otimes C$$ 
$$\lambda : A\lra I\otimes A$$
$$\rho : A\otimes I \lra A$$ 
(A skew monoidal category is {\em Hopf} when $\alpha$ is invertible.)
\item<3-> plus five axioms as follows:

\end{itemize}

\end{frame}


\begin{frame}
  \frametitle{The pentagon and trivial unit condition}
  
\begin{equation}\label{pentagon}
\xymatrix{
& (W\otimes X)\otimes (Y\otimes Z) \ar[rd]^-{{\alpha_{W\otimes X,Y, Z}}}  & \\
W\otimes (X\otimes (Y\otimes Z)) \ar[ru]^-{{\alpha_{W, X,Y\otimes Z}}} \ar[d]_-{1_W\otimes {\alpha_{X, Y,Z}}} & & ((W\otimes X)\otimes Y)\otimes Z  \\
W\otimes ((X\otimes Y)\otimes Z) \ar[rr]_-{{\alpha_{W,X\otimes Y,Z}}} & &  (W\otimes (X\otimes Y))\otimes Z\ar[u]_-{{\alpha_{W,X,Y}} \otimes 1_Z} }
\end{equation}

\begin{equation}\label{unitunit}
\xymatrix{
I \ar[rd]_{\lambda_I}\ar[rr]^{1_I}   && I  \\
& I\otimes I \ar[ru]_{\rho_I} 
}
\end{equation}
\end{frame}

\begin{frame}
  \frametitle{Three more unit conditions}
\begin{equation}\label{leftunit}
\xymatrix{
 I\otimes (X\otimes Y)\ar[rr]^{\alpha_{I,X,Y}}   && (I\otimes X)\otimes Y  \\
& X\otimes Y \ar[lu]^{\lambda_{X\otimes Y}} \ar[ru]_{\lambda_{X} \otimes 1_Y}  &
}
\end{equation}

\begin{equation}\label{midunit}
\xymatrix{
 X\otimes (I\otimes Y)\ar[rr]^-{\alpha_{X,I,Y}} && (X\otimes I)\otimes Y \ar[d]^-{\rho_X \otimes 1_Y} \\
X\otimes Y \ar[u]^-{1_X\otimes \lambda_Y} \ar[rr]_-{1_{X\otimes Y}} && X\otimes Y}
\end{equation}

\begin{equation}\label{rightunit}
\xymatrix{
X\otimes (Y\otimes I) \ar[rr]^{\alpha_{X,Y,I}} \ar[rd]_{1_X\otimes \rho_Y}   && (X\otimes Y)\otimes I  \ar[ld]^{\rho_{X \otimes Y}} \\
& X\otimes Y   &
}
\end{equation}
\end{frame}

\begin{frame}
  \frametitle{Examples from a flock}
  
 \begin{itemize}
\item<1->  Let $\CA$ be a flock. 
\item<2-> Choose any object $J\in \CA$. 
\item<3-> We can define a right skew monoidal structure on $\CA$.
\item<4-> Put $A\otimes B = Q(A,J,B)$ and take $J$ as skew unit.
\item<5-> $\alpha : A\otimes (B\otimes C) \lra (A\otimes B)\otimes C$ is 
$$\phi : Q(A,J,Q(B,J,C)) \lra Q(Q(A,J,B),J,C) \ .$$
\item<6-> $\lambda : A \lra J\otimes A$ is $\zeta : A \lra Q(J,J,A)$.
\item<7-> $\rho : A\otimes J \lra A$ is $\xi : Q(A,J,J) \lra A$.
\end{itemize}

\end{frame}

\begin{frame}
  \frametitle{Example from a bialgebra}
  
 \begin{itemize}
\item<1->  Let $H$ be a bialgebra: so we have linear functions 
$$k \stackrel{\varepsilon}\lla H \stackrel{\delta}\lra H\otimes H \stackrel{\mu}\lra H\stackrel{\eta}\lla k \ .$$ 
\item<2-> The fusion or Galois map $\nu$ is defined to be the composite 
$$\nu : H\otimes H \stackrel{\delta \otimes 1_H}\lra H\otimes H\otimes H  \stackrel{1_H \otimes \mu}\lra H\otimes H \ .$$ 
\item<3-> We can define a skew monoidal structure on $\mathrm{Vect}$ with skew unit $k$, however, with the brand new tensor product $A\star B = H\otimes A \otimes B$.
\item<4-> $\alpha : H\otimes A\otimes H\otimes B\otimes C \lra H\otimes H\otimes A\otimes B\otimes C$ 
does a switch of the second and third factors followed by $\nu \otimes 1_{A\otimes B\otimes C}$.
\item<5-> $\lambda : A \lra H\otimes A$ is $\eta \otimes 1_A$ and $\rho : H\otimes A \lra A$ is $\varepsilon \otimes 1_A$.
\item<6-> This skew monoidal structure on $\mathrm{Vect}$ is Hopf if and only if the bialgebra $H$ is a Hopf algebra.
\end{itemize}

\end{frame}

\begin{frame}
  \frametitle{Example from a category}
  
 \begin{itemize}
\item<1->  Let $A$ be a category. Put $X=\mathrm{ob}A$, assumed small.
\item<2-> $\mathrm{Set}^X$ is the category of families $M=(M_x)_{x\in X}$ of sets. 
\item<3-> We shall define a right-skew monoidal structure on $\mathrm{Set}^X$.
\item<4-> The tensor product $M\ast_A N$ is defined by
$$(M\ast_A N)_x = \sum_y{M_x \times A(x,y) \times N_y} \ .$$
\item<5-> The unit object $J$ is defined by $J_x=1$ for all $x\in X$.
\item<6-> Observe that
$$(M\ast_A (N\ast_A L))_x = \sum_{y,z} {M_x \times A(x,y) \times N_y \times A(y,z)\times L_z}$$
$$((M\ast_A N)\ast_A L)_x = \sum_{y,z} {M_x \times A(x,y) \times N_y \times A(x,z)\times L_z} $$
\end{itemize}
\end{frame}

\begin{frame}
  \frametitle{Example from a category (continued)}
 \begin{itemize}
\item<1->The morphism $\alpha : M\ast_A (N\ast_A L) \lra (M\ast_A N)\ast_A L$ uses the function
$$\nu : A(x,y)\times A(y,z) \lra A(x,y)\times A(x,z)$$
defined by $$\nu(u:x\rightarrow y, v: y\rightarrow z) = (u:x\rightarrow y, v\circ u:x\rightarrow z) \ .$$ 
\item<2-> 
The morphism $\lambda : N \lra J\ast_A N$ has components
$$N_x \stackrel{\mathrm{id_x}\times 1}\lra A(x,x)\times N_x \stackrel{\mathrm{in}_x}\lra \sum_y{A(x,y)\times N_x} \ .$$ 
\item<3-> 
The morphism $\rho : M\ast_A J\lra M$ projects $\sum_y{M_x\times A(x,y) \lra M_x}$. 
\item<4-> This skew monoidal structure on $\mathrm{Set}^X$ is Hopf iff $A$ is a groupoid. 
 \end{itemize}
\end{frame}

 \begin{frame}
  \frametitle{Closed skew monoidal categories}
\begin{itemize}
\item A skew monoidal category $\CR$ is {\em closed} when all the endofunctors 
$A\otimes -$ and $-\otimes A$ on $\CR$ have right adjoints (called {\em internal homs}).
\pause
\item For nice categories like categories of modules, existence of these internal homs
amounts to preservation of colimits by $A\otimes -$ and $-\otimes A$.
 \end{itemize}
  \pause
 \begin{theorem}
Category structures with object set $X$ are equivalent to closed right-skew monoidal structures
on the category $\mathrm{Set}^{X}$ of $X$-indexed families of sets for which the skew unit is the terminal object. 
 \end{theorem}
 \end{frame}
 
\begin{frame}
  \frametitle{Bialgebroids}
\begin{itemize}
\item Under a different name, bialgebroids occurred in the work of Takeuchi (1977). 
They arose again in the quantum group literature. 
\item If you know what bialgebroids are, the following will be a Theorem. Otherwise, it can be a Definition.
  \pause
 \item
 \begin{theorem}[Szlach\'anyi 2012]
Bialgebroids with base ring $R$ are equivalent to closed right-skew monoidal structures
on the additive category $\mathrm{Ab}^{R^{\circ}}$ of right $R$-modules for which the skew unit is $R$. 
 \end{theorem}
 \pause
 \item The proof begins by noting that tensor products 
 $$\mathrm{Ab}^{R^{\circ}} \otimes \mathrm{Ab}^{R^{\circ}} \lra \mathrm{Ab}^{R^{\circ}} \ ,$$
 which are colimit-preserving in each variable, are induced by tensoring with a left 
 $R\otimes R$-, right-$R$-bimodule $A$. This $A$ becomes a ring with two
 ring morphisms $s : R^{\circ} \lra A$ and $t : R \lra A$ with commuting images. 
\end{itemize}

\end{frame}
 \begin{frame}
  \frametitle{Two more skew monoidal categories}
  \begin{itemize}
\item Write $\mathbf{Fsk}$ for the free right skew monoidal category on a single generating object $N$.
\pause 
\item For each object $A$ of any right skew monoidal category $\CR$, there exists a unique
strict monoidal functor $\mathbf{Fsk} \lra \CR$ taking $N$ to $A$. 
\pause
\item In the joint work with Steve Lack, we have a concrete combinatorial model for this.
\pause
\item Let $\mathbf{\Delta}_{\top}$ denote the category whose objects are the finite ordinals 
$\mathbf{n} = \{0,1,\dots , n-1\}$
with last element (so that $n>0$), 
and whose morphisms are last-element-and-order-preserving functions.
\pause
\item $\mathbf{\Delta}_{\top}$ is Hopf left skew monoidal with ordinal sum as
tensor product, with unit $\mathbf{1}$, with right unit constraint $\rho : \mathbf{n+1}\lra \mathbf{n}$
the surjection identifying $n-1$ and $n$, and with left unit constraint $\lambda : \mathbf{m}\lra \mathbf{1+m}$ the injection whose image does not contain $0$.   
\end{itemize}
\end{frame}

 \begin{frame}
  \frametitle{Coherence for skew monoidal categories}
  \begin{itemize}
\item Monoidal categories are skew monoidal categories in which all the constraints are invertible. 
\pause
\item Mac Lane (1963) proved a `coherence theorem' for monoidal categories: roughly, 
``all diagrams made up of constraints and their inverses commute''.  
\pause
\item More precisely, all diagrams in the free monoidal category $\mathbf{Fmc}$ on a single generating object commute. In other words, the unique functor $\mathbf{Fmc} \lra \mathbf{1}$ is faithful. 
\end{itemize}
\pause
\begin{theorem}[Lack-Street 2013]
The unique strict monoidal functor $\mathbf{Fsk} \lra \mathbf{\Delta}_{\top}$, which takes the 
generating object $N$ of $\mathbf{Fsk}$ to the object $\mathbf{1}$ of 
$\mathbf{\Delta}_{\top}$, is faithful.  
\end{theorem}
\end{frame}
\end{document}
